% Definizione di un verde leggibile come nelle slide
\definecolor{slidegreen}{RGB}{0, 166, 81}

\chapter{2ª Esercitazione}

\section{Alternanza pari e dispari}
\subsection{Definizione del problema}

\subsubsection{Istanza}
Una sequenza $X = \langle x_1,x_2, \dots, x_n \rangle$ di numeri

\subsubsection{Soluzione}
Una sottosequenza $Z = \langle z_1, z_2, \dots, z_k \rangle$ tale che:
\begin{itemize}
    \item In $Z$ si alternino elementi pari e dispari
    \item $k$ sia massimo
\end{itemize}

\subsection{Definizione del sottoproblema}
Travare la più lunga sottosequenza $ALT_{V}(X_i)$ con $i \in \{1,\dots,n\}$ del prefisso $X_i$, che alterna pari e dispari che termina con l'elemento $x_i$

\subsection{Definizione del coefficiente $\mathbf{c_i}$}
$c_i$ è la lunghezza della sequenza $ALT1_V(X_i)$

\subsection{Formulazione ricorsiva}

\subsubsection{Caso base}
\begin{itemize}
    \item $\mathbf{i=1} \quad X_1= \langle x_1 \rangle$ \\
    $\implies ALT1_V(X_1) = \langle x_1 \rangle$ \\
    $\implies c_1 = 1$
\end{itemize}

\subsubsection{Passo ricorsivo}
\begin{itemize}
    \item $\mathbf{i > 1}$ \\
    $\implies ALT_V(X_i) = longest\{ALT1_V(X_j) : 1 \leq j < i, x\%2 \not = x_j \%2\} + \langle x_i \rangle$ \\
    $\implies c_i = \max\{c_j : 1 \leq j < i, x_i\%2 \not = x_j \% 2 \} + 1$
\end{itemize}

\newpage
\section{Alternanza dei colori}

\subsection{Definizione del problema}
\subsubsection{Istanza}
\begin{itemize}
    \item $X = \langle x_1, x_2, \dots, x_n \rangle$
    \item $col:\Sigma \to \{red, green, black\}$
\end{itemize}

\subsubsection{Soluzione}
Una sottosequenza $Z = \langle z_1, z_2, \dots, z_k \rangle$ tale che:
\begin{itemize}
    \item in $Z$ si alternino elementi rossi e neri
    \item $k$ sia massimo
\end{itemize}

\subsection{Definizione del sottoproblema}
Trovare la più lunga sottosequenza $ALT_V(X_i)$ del prefisso $X_i$ con $i \in \{1,\dots,n\}$ che alternino rossi e neri che termina con l'elemento $x_i$

\subsection{Formulazione ricorsiva}

\subsubsection{Caso base}
\begin{itemize}
    \item $i=1 \land col(x_1) \not = green$ \\
    $\implies ALT_V(X_1) = \langle x_1 \rangle$ \\
    $\implies c_1 = 1$

    \item $\forall i, col(x_i) = green$ \\
    $\implies ALT_V(X_i) = NIL$ \\
    $\implies c_i = 0$
\end{itemize}

\subsubsection{Passo ricorsivo}
\begin{itemize}
    \item $i>1 \land col(x_i) \not = green$\\
    $\implies ALT_V(X_i) = longest\{ALT_V(X_i) \not = NIL : 1 \leq j < i, col(x_i) \not = col(x_j)\} + \langle x_i \rangle$ \\
    $\implies c_i = \max\{c_j : 1 \leq j < i, col(x_j) \not = col(x_i)\} +1$
\end{itemize}

\newpage
\section{Elementi non consecutivamenti rossi}

\subsection{Definizione del problema}
\subsubsection{Istanza}
\begin{itemize}
    \item $X = \langle x_1, x_2, \dots, x_n \rangle$
    \item $col: \Sigma \to \{red, green, black\}$
\end{itemize}

\subsubsection{Soluzione}
Una sottosequenza $Z = \langle z_1, z_2, \dots, z_k \rangle$ tale che:
\begin{itemize}
    \item In $Z$ non vi sono elementi consecutivi rossi
    \item $k$ sia massimo
\end{itemize}

\subsection{Definizione del sottoproblema}
Trovare la più lunga sottosequenza $NC_V(X_i)$ del prefisso $X_i$ in cui non vi sono elementi consecutivi rossi e che termina con l'elemento $x_i$

\subsection{Definizione del coefficiente}
Lunghezza di $NC_V(X_i)$

\subsection{Formulazione ricorsiva}

\subsubsection{Caso base}
\begin{itemize}
    \item $\mathbf{i=1}$ \\
    $\implies NC_V(X_i) = \langle x_i \rangle$
    $\implies c_i = 1$
\end{itemize}

\subsubsection{Passo ricorsivo}
\begin{itemize}
    \item $i>1$\\
    $\implies NC_V(X_i) = longest\{NC_V(X_j): 1 \leq j < i, col(x_j) \not = red \lor col(x_i) = red\} + \langle x_i \rangle $ \\
    $\implies c_i = \max\{c_j : 1 \leq j < i, col(x_j) \not = red \lor col(x_i) \not = red\} + \langle x_i \rangle$
\end{itemize}

\newpage
\section{Elementi non consecutivamente dello stesso colore}

\subsection{Definizione del problema}
\subsubsection{Istanza}
\begin{itemize}
    \item $X = \langle x_1, x_2, \dots, x_n \rangle$
    \item $col: \Sigma \to \{red, green, black\}$
\end{itemize}

\subsubsection{Soluzione}
Una sottosequenza $Z = \langle z_1, z_2, \dots, z_k \rangle$ tale che:
\begin{itemize}
    \item in $Z$ non vi siano elementi consecutivi dello stesso colore
    \item $k$ sia massimo
\end{itemize}

\subsection{Definizione del sottoproblema}
Trovare la più lunga sottosequenza $NC_V(X_i)$ del prefisso $X_i$, in cui non vi siano elementi consecutivi dello stesso coloro e che termina con l'elemento $x_i$

\subsection{Definizione del coefficiente}
$c_i$ è la lunghezza di $NC_V(X_i)$


\subsection{Formulazione ricorsiva}

\subsubsection{Caso base}
\begin{itemize}
    \item $i = 1$ \\
    $\implies NC_V(X_1) = \langle x_1 \rangle$ \\
    $\implies c_1 = 1$
\end{itemize}


\subsubsection{Passo ricorsivo}
\begin{itemize}
    \item $i>1$ \\
    $\implies NC_V(X_i) = longest\{NC_V(x_j) : 1 \leq j < i, col(x_j) \not = col(x_i)\} + \langle x_i \rangle$ \\
    $\implies c_i = \max\{c_j : 1 \leq j < i, col(x_j) \not = col(x_i)\} + 1$
\end{itemize} 

\newpage
\section{Problema della Edit Distance}
\subsection{Testo del problema}

$S(X,Y)$ è la minima sequenza di operazioni di cancellazione, inserimento, sostituzione di un singolo elemento da applicare su $X$ (da sinistra) per trasformarla in Y

\begin{enumerate}
    % Passo 1
    \item \textbf{Sequenze iniziali:}
    \[
    \begin{aligned}
    X &= \langle 2, 4, 10, 3, 1 \rangle \\
    Y &= \langle 2, 3, 4, 2, 1 \rangle
    \end{aligned}
    \]

    % Passo 2
    \item \textbf{Inserimento di $\textcolor{slidegreen}{\mathbf{3}}$:}
    \[
    \begin{aligned}
    X &= \langle 2, \underset{\textcolor{slidegreen}{\vee}}{\phantom{,}} 4, 10, 3, 1 \rangle \\
    Y &= \langle 2, 3, 4, 2, 1 \rangle
    \end{aligned}
    \]
    \[
    S(X, Y) = \langle \text{inserisco } \textcolor{slidegreen}{\mathbf{3}}
    \]

    % Passo 3
    \item \textbf{Cancellazione di $\textcolor{red}{\mathbf{10}}$:}
    \[
    \begin{aligned}
    X &= \langle 2, \overset{\textcolor{slidegreen}{\mathbf{3}}}{\underset{\textcolor{slidegreen}{\vee}}{\phantom{,}}}, 4, \textcolor{red}{\xcancel{\mathbf{10}}}, 3, 1 \rangle \\
    Y &= \langle 2, 3, 4, 2, 1 \rangle
    \end{aligned}
    \]
    \[
    S(X, Y) = \langle \text{inserisco } \textcolor{slidegreen}{\mathbf{3}}, \text{ cancello } \textcolor{red}{\mathbf{10}}
    \]

    % Passo 4 e 5
    \item \textbf{Sostituzione di $\textcolor{blue}{\mathbf{3}}$ con $2$ e calcolo di $ED(X, Y)$:}
    \[
    \begin{aligned}
    X &= \langle 2, \overset{\textcolor{slidegreen}{\mathbf{3}}}{\underset{\textcolor{slidegreen}{\vee}}{\phantom{,}}}, 4, \textcolor{red}{\xcancel{\mathbf{10}}}, \overset{\textcolor{blue}{\mathbf{2}}}{\textcolor{blue}{\cancel{\mathbf{3}}}}, 1 \rangle \\
    Y &= \langle 2, 3, 4, 2, 1 \rangle
    \end{aligned}
    \]
    \begin{align*}
    S(X, Y) &= \langle \text{inserisco } \textcolor{slidegreen}{\mathbf{3}}, \text{ cancello } \textcolor{red}{\mathbf{10}}, \text{ sostituisco } \textcolor{blue}{\mathbf{3}} \text{ con } 2 \rangle \\
    ED(X, Y) &= |S(X, Y)| = 3 \qquad (\text{Edit Distance})
    \end{align*}

    % Passo 6 e 7
    \item \textbf{Sequenza alternativa $S(X, Y)$:}
    \[
    \begin{aligned}
    X &= \langle 2, \overset{\textcolor{slidegreen}{\mathbf{3}}}{\underset{\textcolor{slidegreen}{\vee}}{\phantom{,}}}, 4, \overset{\textcolor{red}{\mathbf{2}}}{\textcolor{red}{\cancel{\mathbf{10}}}}, \textcolor{blue}{\xcancel{\mathbf{3}}}, 1 \rangle \\
    Y &= \langle 2, 3, 4, 2, 1 \rangle
    \end{aligned}
    \qquad \text{Un'altra sequenza } S(X, Y)\text{?}
    \]
    \begin{align*}
    S(X, Y) &= \langle \text{inserisco } \textcolor{slidegreen}{\mathbf{3}}, \text{ sostituisco } \textcolor{red}{\mathbf{10}} \text{ con } 2, \text{ cancello } \textcolor{blue}{\mathbf{3}} \rangle \\
    ED(X, Y) &= |S(X, Y)| = 3 \qquad (\text{Edit Distance})
    \end{align*}
\end{enumerate}

\subsection{Definizione del problema}
\subsubsection{Istanza}
\begin{itemize}
    \item $X = \langle x_1, \dots, x_m \rangle$ 
    \item $Y = \langle y_1, \dots, y_n \rangle$
\end{itemize}

\subsubsection{Soluzione}
Una sequenza $S(X,Y) = \langle o_1, \dots, o_k \rangle$ di operazioni di cancellazione, inserimento e sostituzione che trasformano $X$ in $Y$ tale che $k$ sia minimo

\subsection{Definizione del sottoproblema}
Trovare $S(X_i,Y_j)$ dove:
\begin{itemize}
    \item $i \in \{0, \dots, m\}$
    \item $j \in \{0, \dots, n\}$
\end{itemize}

\subsection{Definizione del coefficiente}
$c_{ij}$ è la lunghezza di $S(X_i, Y_j)$

\subsection{Formulazione ricorsiva}
\subsubsection{Caso base}
\begin{itemize}
    \item $i = 0 \land j=0$\\
    $\implies S(X_i, Y_j) = \langle \rangle$\\
    $c_{ij} = 0$
\end{itemize}

\subsubsection{Passo ricorsivo}
\begin{itemize}
    \item $(i > 0 \land j > 0) \land (x_i \not = y_j)$ \\
    $\implies S(X_i, Y_j) = shortest 
    \begin{cases}
        S(X_{i-1}, Y_{j-1}) + \langle \text{subst } x_i \text{ con } y_i \rangle \\
        S(X_{i-1},Y_j) + \langle \text{delete } x_i\rangle \\
        S(X_i, Y_{j-1}) + \langle \text{insert } y_j \rangle
    \end{cases}
    $ \\
    $\implies c_{ij} = min 
    \begin{cases}
        c_{i-1\ j-1} + 1 \\
        c_{i-1\ j} + 1 \\
        c_{i\ j-1} + 1    
    \end{cases}
    $
\end{itemize}